\section{Soft Robotics and Pneumatic Actuation}
\label{sec:soft-robotics}

The simplest way to ask what makes a soft robot soft is to look at what it cannot do. A rigid manipulator with revolute joints can place an end-effector at a precise pose, but it cannot bend around an unexpected obstacle, conform to an irregular surface, or absorb a collision without damage. The reason is constitutive: the material it is made of permits no large strain, so all motion must be accumulated at discrete joints, and all interaction with the world must be specified by the controller. Soft robotics begins with a different premise. Its bodies are made of materials such as silicone elastomers, fluidic networks, and inflatable membranes, whose Young's modulus lies orders of magnitude closer to biological tissue than to engineering steel \cite{rusDesignFabricationControl2015}, and which deform continuously under load, distributing strain across the whole body rather than concentrating it at a hinge. The shift in material is also a shift in what the controller has to do, since deformation the body absorbs on its own is deformation the controller no longer has to specify. Reviews of the field frame this trade-off as the central promise of soft robotics: that compliance enables abilities such as squeezing, stretching, climbing, growing, and morphing, which are unreachable with rigid links alone \cite{laschi_soft_2016,rusDesignFabricationControl2015}.

This material-level shift carries an actuation-level consequence that splits the control problem in two: \textit{actuation becomes simpler, while modeling becomes harder}. On the actuation side, soft bodies absorb much of what a rigid controller would need to specify. As Wang and colleagues put it, \inlinequote{the inherent material softness can facilitate and simplify control}\cite{wangEmbodyingPhysicalComputing2026}; a single global pressure input is often enough to make a soft tentacle wrap an object or a soft leg accommodate uneven terrain. On the modeling side, the same softness defeats the analytical tools developed for rigid robotics: \inlinequote{The kinematics and dynamics of soft robotic systems are unlike those of traditional, rigid-bodied robots.}\cite{rusDesignFabricationControl2015} The continuous deformation of a soft body cannot be reduced to the joint angles of a kinematic chain, and the inverse mapping from actuation to motion is strongly nonlinear, continuous, and underdetermined.

Pneumatic actuation is one of the most direct illustrations of this trade-off. Pressurizing an enclosed elastic chamber changes its volume and, with appropriate geometric constraints, its shape; energy flows through the material rather than across a kinematic chain. The control input is sparse, since a single pump and a small set of valves can drive many inflatable elements, decoupling actuator count from module count when the design must scale. What each chamber does in response to a given pressure, however, depends on its geometry, its material, the loading from neighbors, and the contact forces from the environment, none of which are cleanly captured by classical rigid-body models. The soft-robotics trade-off, in short, is to hand part of the control problem over to the body and the material, accepting in exchange that the part that remains is harder to model than before.

Three systems make this trade-off tangible. Kim, Kim, and Park \cite{kimSimpleTripodMobile2019} build a tripod robot in which each leg is an inflatable membrane coupled to a passive valve shaft. A constant input pressure inflates the membrane until expansion physically opens an exhaust port; the membrane then deflates, the valve closes, and the cycle repeats. The frequency of this self-oscillating cycle is set by the steady input pressure, not by a time-varying control signal: the controller specifies a value, and the gait is what the membrane does in response. Wait, Jackson, and Smoot \cite{waitSelfLocomotionSpherical2010} take a different route. Their spherical robot has an outer surface composed of thirty-two individually inflatable rubber bladders, arranged on the faces of a truncated icosahedron (\textit{\cref{fig:wait_spherical}}). Inflating bladders in the region near the ground, but offset from the desired direction of travel, deforms the surface at the contact point and applies a rolling moment to the sphere; locomotion here is the geometric consequence of \textit{which} bladders are inflated, not a sequence of joint commands. Mousa et al. \cite{mousaMultifunctionalFluidicUnits} push the idea further still. Their multifunctional fluidic units, each capable of behaving simultaneously as valve, sensor, and actuator, can be assembled into networks that produce coordinated locomotion from minimal input. More striking, when several such oscillating units are mechanically coupled through a shared elastic body, they exhibit emergent self-synchronization, modeled as a Kuramoto network of coupled oscillators. The coordination between actuators is not specified anywhere in the controller; it arises passively, through the mechanical coupling of the units themselves. Across the three systems, the same pattern recurs: a property the controller would otherwise need to specify, such as gait timing, body posture, or inter-actuator coordination, is instead absorbed by the material's own response to a small pneumatic input.

What these systems demonstrate, taken together, is that soft robotics is a way of letting material and form take on part of the computation that a controller would otherwise have to perform. This is the territory of morphological computation, which the next section examines in its own right.

The soft-robotics literature also leaves open at least two questions that resonate with the dialectical pairs of Section~\ref{sec:dialectics}. The modeling and control difficulties that follow from a fully soft body raise the question of whether soft and rigid components can be deliberately combined: a body composed of compliant elements held together by a discrete rigid skeleton might preserve the body-as-computation property of the soft parts while keeping the rigid parts tractable to analyze. The soft-versus-rigid pair returns here as a concrete design choice rather than a metaphysical one. In a similar spirit, the strong nonlinearity and underdetermination that make soft robots hard to model raise the question of whether part of that unpredictability is a resource rather than a defect. The predictable-versus-unpredictable pair from the same section suggested that physical systems compute partly because they are not fully reducible to symbolic specification; taken at the level of body design, this implies that the unmodelable parts of a soft robot's behavior need not be treated as errors to be filtered out, but as candidate contributions to the robot's overall computation.

Each of the three systems above is a single, monolithic body, or in Mousa's case a single planar fluidic network, whose form was chosen by its designer and fixed thereafter; comparison across designs is qualitative at best. Laschi, Mazzolai, and Cianchetti note in passing that \inlinequote{inflatable robots are based on multiple inflatable chambers whose shape determines the robot motion}\cite{laschi_soft_2016}. The shape, in other words, is doing some of the work, but the question of \textit{which} shape, and \textit{why}, has been left almost untouched. That is the question the next section takes up.
